% =============================================================================
%  claude_atm_2.tex
%
%  A proof of \Cref{thm:main} (the unramifiedness theorem of
%  proof_complete_standalone.tex) routed through *local* ramification theory:
%  the cyclic-Herbrand convexity lemma is the new ingredient, and a
%  separately-cited "local--global dictionary" relating geometric
%  unramifiedness of $\cP^{[d]_\Gamma}$ at $\eta_C$ to the upper-ramification
%  bound $\Gamma^d=1$ at $p_\infty$ does the rest.
%
%  This proof is *for comparison only* with the direct Cartesian-inertia
%  argument of claude_atm_1.tex. It is strictly less self-contained --- the
%  load-bearing claim is the dictionary, not the new lemma --- and we are
%  explicit about that.
%
%  Compile with:  latexmk -pdf claude_atm_2.tex   (uses biber)
% =============================================================================
\documentclass[11pt]{article}
\usepackage{geometry}\geometry{margin=1in}
\usepackage[T1]{fontenc}
\usepackage{ifthen}

\setlength{\parindent}{1em}
\setlength{\parskip}{0.5em}

\input{includes}
\input{MacrosAndFigures}

\usepackage{csquotes}
\usepackage[
backend=biber,
style=alphabetic,
sorting=ynt
]{biblatex}
\addbibresource{Bib.bib}

\newboolean{ThesisIntro}
\setboolean{ThesisIntro}{false}

\title{Unramifiedness of $\cP^{[d]_G}$ at the exceptional divisor\\[2pt]
       {\large (Herbrand-convexity route, for comparison)}}
\author{Assaf Marzan}
\date{\today}

\begin{document}
\maketitle

\noindent
\textbf{Scope.} This note gives a proof of \Cref{thm:main-herb} (= \texttt{thm:main}
of \texttt{proof\_complete\_standalone.tex}) routed through local ramification
theory. The new content is a convexity lemma for the Herbrand $\varphi$-function
on a cyclic $p$-group (\Cref{lem:herbrand-convexity}). The route works only if
one accepts the local--global dictionary stated in \Cref{prop:dictionary} below
as a separately-established result; we discuss the resulting circularity at the
end (\Cref{rem:circularity}).

We adopt the notation of \texttt{proof\_complete\_standalone.tex} without
restating the background: $C$ is a smooth projective curve over $k$, $p_\infty
\in C$, $U = C \setminus \{p_\infty,\ldots\}$, $\cP \to U$ a $G$-torsor with
$G \cong \Z/p^r\Z$ a cyclic $p$-group, $H \subseteq G$ a subgroup of order
$p^s$, $f \colon C' \to C$ the compactification of the $G/H$-torsor
$\cP_{G/H} \to U$, $p'_\infty \in f^{-1}(p_\infty)$, and $X_C, X_{C'}$ the
blowups of $C^{(d)}, C'^{(d)}$ at $d\cdot p_\infty, d \cdot p'_\infty$ with
exceptional generic points $\eta_C, \eta_{C'}$. Throughout, $0 < s < r$
(the case $s = 0$ or $s = r$ being trivial).

% #############################################################################
\section{The Herbrand convexity lemma}\label{sec:herbrand}
% #############################################################################

The single new local ingredient. We work in the setting of a finite totally
ramified Galois extension of complete (equivalently, henselian and defectless)
discrete valuation fields with cyclic $p$-group Galois group. Lower and upper
numbering for ramification follow \cite{Serre1979}, Ch.~IV.

\begin{lemma}\label{lem:herbrand-convexity}
Let $\mathcal{F}/\mathcal{K}$ be a finite totally ramified Galois extension of
complete discrete valuation fields with Galois group $\Gamma \cong \Z/p^r\Z$,
$r \geq 1$, and let $\Delta \subseteq \Gamma$ be the unique subgroup of order
$p^s$ with $0 < s < r$, so that $\Gamma/\Delta \cong \Z/p^{r-s}\Z$ and
$|\Delta| = p^s$. Fix an integer $d \geq 1$. If both
\begin{enumerate}
  \item[\textnormal{(a)}] $(\Gamma/\Delta)^d = \{1\}$, i.e.\ the upper
        ramification filtration of $\Gamma/\Delta$ acting on
        $\mathcal{F}^\Delta/\mathcal{K}$ is trivial in degree $\geq d$, and
  \item[\textnormal{(b)}] $\Delta^d = \{1\}$, i.e.\ the upper ramification
        filtration of $\Delta$ acting on $\mathcal{F}/\mathcal{F}^\Delta$ is
        trivial in degree $\geq d$,
\end{enumerate}
then $\Gamma^d = \{1\}$, i.e.\ the upper ramification filtration of $\Gamma$
acting on $\mathcal{F}/\mathcal{K}$ is trivial in degree $\geq d$.
\end{lemma}

\begin{proof}
Let
\[
\Gamma = \Gamma_0 \supsetneq \Gamma_{m_1+1} \supsetneq \Gamma_{m_2+1} \supsetneq
\cdots \supsetneq \Gamma_{m_r+1} = \{1\}
\]
be the lower-numbering ramification filtration of $\Gamma$ on
$\mathcal{F}/\mathcal{K}$, with jumps at integers $m_1 < m_2 < \cdots < m_r$.
Since $|\Gamma| = p^r$ and the extension is totally ramified, there are exactly
$r$ jumps, and $|\Gamma_{m_j+1}| = p^{r-j}$. Set
\[
\Delta_0 := m_1, \qquad \Delta_j := m_{j+1} - m_j \quad (1 \leq j \leq r-1),
\]
so that $m_j = \Delta_0 + \Delta_1 + \cdots + \Delta_{j-1}$. (We reuse the
letter $\Delta$ for the gap sequence; the subgroup $\Delta \subseteq \Gamma$
will appear only as such, with subscripts to differentiate.)

\paragraph{The three Herbrand functions.}
Recall $\varphi_\Gamma(u) = \int_0^u |\Gamma_t|/|\Gamma_0|\,dt$. The integrand
takes the constant value $p^{-j}$ on $(m_j, m_{j+1}]$ ($0 \leq j \leq r-1$,
with $m_0 := 0$), so
\begin{equation}\label{eq:phi-Gamma}
\varphi_\Gamma(m_r) \;=\; \sum_{j=0}^{r-1} p^{-j}\,\Delta_j.
\end{equation}

For the subgroup $\Delta = \Gamma_{m_{r-s}+1}$ acting on
$\mathcal{F}/\mathcal{F}^\Delta$, the lower ramification of $\Delta$ as a
subgroup of $\Gamma$ is given by $\Delta_i = \Delta \cap \Gamma_i$
(\cite{Serre1979}, Ch.~IV, \S1, Prop.~2). Therefore $\Delta_i = \Delta$ for
$0 \leq i \leq m_{r-s}$, $\Delta_i = \Gamma_i$ for $i > m_{r-s}$, and the lower
jumps of $\Delta$ occur precisely at $m_{r-s+1} < m_{r-s+2} < \cdots < m_r$.
Hence
\begin{equation}\label{eq:phi-Delta}
\varphi_\Delta(m_r) \;=\; m_{r-s+1}
   + \sum_{k=1}^{s-1} p^{-k}(m_{r-s+k+1} - m_{r-s+k})
   \;=\; m_{r-s} + \sum_{k=0}^{s-1} p^{-k}\,\Delta_{r-s+k}.
\end{equation}

For the quotient $\Gamma/\Delta$ acting on $\mathcal{F}^\Delta/\mathcal{K}$,
Herbrand's theorem (\cite{Serre1979}, Ch.~IV, \S3, Lemma~5) gives
$(\Gamma/\Delta)^v = \Gamma^v \Delta / \Delta$. Writing
$u_j := \varphi_\Gamma(m_j)$ for the upper jumps of $\Gamma$, one has
$\Gamma^{u_j} = \Gamma_{m_j}$, so $\Gamma^{u_j} \Delta \supsetneq \Delta$ iff
$\Gamma_{m_j} \not\subseteq \Delta$, iff $j \leq r-s$. The highest upper jump of
$\Gamma/\Delta$ is therefore $u_{r-s} = \varphi_\Gamma(m_{r-s})$, and
\begin{equation}\label{eq:phi-quotient}
u_{r-s} \;=\; \varphi_\Gamma(m_{r-s}) \;=\; \sum_{j=0}^{r-s-1} p^{-j}\,\Delta_j.
\end{equation}

\paragraph{The fundamental identity.}
Split the sum in \eqref{eq:phi-Gamma} at index $j = r-s$ and factor
$p^{-(r-s)}$ from the tail:
\begin{align*}
\varphi_\Gamma(m_r)
 &= \underbrace{\sum_{j=0}^{r-s-1} p^{-j}\Delta_j}_{= u_{r-s}\text{ by }\eqref{eq:phi-quotient}}
   \;+\; p^{-(r-s)}\underbrace{\sum_{k=0}^{s-1} p^{-k}\Delta_{r-s+k}}_{= \varphi_\Delta(m_r) - m_{r-s}\text{ by }\eqref{eq:phi-Delta}}.
\end{align*}
Hence
\begin{equation}\label{eq:fund-identity}
\boxed{\;\varphi_\Gamma(m_r) \;=\; u_{r-s} \;+\; \frac{\varphi_\Delta(m_r) - m_{r-s}}{p^{r-s}}\;.\;}
\end{equation}

\paragraph{Convex combination.}
Set $A := u_{r-s}$ and $B := \varphi_\Delta(m_r)$. Hypotheses (a) and (b) read
$A < d$ and $B < d$ respectively, by the characterisation of the highest upper
jump as the smallest $v$ such that the upper group is trivial in degree $> v$.
Since the integrand defining $\varphi_\Gamma$ is bounded above by $1$, we have
$\varphi_\Gamma(m_{r-s}) \leq m_{r-s}$, i.e.\ $A \leq m_{r-s}$. Substituting
into \eqref{eq:fund-identity}, $B - m_{r-s} \leq B - A$, hence
\[
\varphi_\Gamma(m_r) \;\leq\; A + \frac{B - A}{p^{r-s}}
                    \;=\; \left(1 - \tfrac{1}{p^{r-s}}\right) A
                          + \tfrac{1}{p^{r-s}}\, B.
\]
The right-hand side is a convex combination of two quantities both strictly
less than $d$, hence
\[
\varphi_\Gamma(m_r) \;<\; \left(1 - \tfrac{1}{p^{r-s}}\right) d
                          + \tfrac{1}{p^{r-s}}\, d \;=\; d.
\]
The highest upper jump of $\Gamma$ is $\varphi_\Gamma(m_r) < d$, so
$\Gamma^d = \{1\}$, as claimed.
\end{proof}

\begin{remark}
The cyclic prime-power hypothesis is essential: it forces both the chain of
lower-numbering ramification groups to pass through every subgroup of $\Gamma$
(which lets us identify the lower jumps of $\Delta$ as a tail of the jumps of
$\Gamma$) and the existence of a unique $\Delta$ of each order (so the
hypotheses (a), (b) make unambiguous sense). For a non-cyclic abelian
$p$-group, the analogous statement fails; the convexity step still works but
the lower jumps of $\Delta \subseteq \Gamma$ need not be a tail.
\end{remark}

% #############################################################################
\section{The dictionary, as a black box}\label{sec:dictionary}
% #############################################################################

The geometric-to-local translation we need takes the following form.

\begin{proposition}[Local--global dictionary for symmetric-power torsors]\label{prop:dictionary}
Let $Z$ be a smooth projective geometrically connected curve over a field $k$,
$q \in Z$ a closed point, $\Gamma$ a finite abelian group, and $\cQ \to Z
\setminus \{q\}$ a $\Gamma$-torsor with associated complete local extension
$\mathcal{F}/\mathcal{K}_q$ at $q$ (in the sense of
\Cref{definition:local_ramification_boundness} of
\texttt{proof\_complete\_standalone.tex}). Let $\pi_Z \colon X_Z \to Z^{(d)}$
be the blowup at $d\cdot q$, with exceptional divisor of generic point
$\eta_Z$. Then
\[
\cQ^{[d]_\Gamma} \to (Z\setminus\{q\})^{(d)}\text{ is unramified at } \eta_Z
\quad\Longleftrightarrow\quad \Gamma^d = \{1\}\text{ on }
\mathcal{F}/\mathcal{K}_q .
\]
\end{proposition}

We will use this proposition three times in the proof of \Cref{thm:main-herb}.
We discuss its status in \Cref{rem:circularity}.

% #############################################################################
\section{The proof}\label{sec:proof}
% #############################################################################

\begin{theorem}\label{thm:main-herb}
Suppose
\begin{enumerate}
  \item[\textnormal{(H1)}] $(\cP_{G/H})^{[d]} \to U^{(d)}$ is unramified at
        $\eta_C$;
  \item[\textnormal{(H2)}] $\cP^{[d]_H} \to (\cP_{G/H})^{(d)}$ is unramified at
        $\eta_{C'}$;
  \item[\textnormal{(H3)}] $G = \Z/p^r\Z$ is cyclic of prime-power order.
\end{enumerate}
Then $\cP^{[d]_G} \to U^{(d)}$ is unramified at $\eta_C$.
\end{theorem}

\begin{proof}
After the standard reductions (replacing $G$ by the image of monodromy and $H$
by its intersection; we may assume $\cP$ is a connected $G$-torsor and that
$0 < |H| < |G|$), we work with the complete local extension
$\mathcal{F}/\widehat{\mathcal{K}}_{p_\infty}$ corresponding to $\cP \to U$ at
$p_\infty$, which is totally ramified with Galois group $G$ (after passing to a
connected component of $\cP \mid_{\Spec \widehat{\mathcal{K}}_{p_\infty}}$ and
extending scalars if needed). The subextension $\mathcal{F}^H \subseteq
\mathcal{F}$ has $\Gal(\mathcal{F}^H/\widehat{\mathcal{K}}_{p_\infty}) = G/H$
and $\Gal(\mathcal{F}/\mathcal{F}^H) = H$. The choice of $p'_\infty \in
f^{-1}(p_\infty)$ corresponds to choosing a place of $\mathcal{F}^H$, with
completion identified with $\widehat{\mathcal{K}}_{p'_\infty}$; the local
extension of the $H$-torsor $\cP \to \cP_{G/H}$ at $p'_\infty$ is then
$\mathcal{F}/\mathcal{F}^H$.

\paragraph{Step 1: From (H1) to a local jump bound for $G/H$.}
Apply \Cref{prop:dictionary} to the $(G/H)$-torsor $\cP_{G/H} \to U$ at
$p_\infty$, with $\Gamma = G/H$. Hypothesis (H1) is the geometric side of the
equivalence, so the local side gives
\[
(G/H)^d = \{1\} \text{ on } \mathcal{F}^H/\widehat{\mathcal{K}}_{p_\infty}.
\]

\paragraph{Step 2: From (H2) to a local jump bound for $H$.}
Apply \Cref{prop:dictionary} to the $H$-torsor $\cP \to \cP_{G/H}$ at the
closed point $p'_\infty \in C'$, with base curve $Z = C'$, group $\Gamma = H$,
and the same integer $d$. The blowup of $C'^{(d)}$ at $d \cdot p'_\infty$ is
exactly $X_{C'}$ of our master diagram, with exceptional generic point
$\eta_{C'}$; the geometric side of the equivalence is therefore (H2). The local
side reads
\[
H^d = \{1\} \text{ on } \mathcal{F}/\mathcal{F}^H.
\]

\paragraph{Step 3: Apply the Herbrand convexity lemma.}
\Cref{lem:herbrand-convexity} applies with $\mathcal{K} =
\widehat{\mathcal{K}}_{p_\infty}$, $\mathcal{F} = \mathcal{F}$, $\Gamma = G$,
$\Delta = H$: hypotheses (a) and (b) are Steps~1 and~2 respectively, and (H3)
is the cyclic prime-power assumption. We conclude
\[
G^d = \{1\} \text{ on } \mathcal{F}/\widehat{\mathcal{K}}_{p_\infty}.
\]

\paragraph{Step 4: Translate back via the dictionary.}
Apply \Cref{prop:dictionary} once more, now to the $G$-torsor $\cP \to U$ at
$p_\infty$ with $\Gamma = G$. The local side is Step~3; the geometric side
reads
\[
\cP^{[d]_G} \to U^{(d)} \text{ is unramified at } \eta_C,
\]
which is the conclusion of \Cref{thm:main-herb}. \qedhere
\end{proof}

% #############################################################################
\section{Discussion: where this proof is circular, and how to rescue it}\label{sec:discussion}
% #############################################################################

\begin{remark}[Circularity in Step~4]\label{rem:circularity}
The application of \Cref{prop:dictionary} in Step~4 of the proof above is for
the group $G$ itself --- the same group that appears in the conclusion of
\Cref{thm:main-herb}. If \Cref{prop:dictionary} is read as a known theorem
established elsewhere (say in the cited literature for sufficiently general
abelian $\Gamma$), this is fine; but if the dictionary is itself the subject of
the work and \Cref{thm:main-herb} is one of its building blocks, then Step~4 is
circular.

The standard way to rescue the argument is to embed it in an \emph{induction on
$|G|$}. The dictionary for cyclic $p$-groups of order $< p^r$ is then the
inductive hypothesis, Steps~1 and~2 use the dictionary in the easy direction
for the strictly smaller groups $G/H$ and $H$, and Step~4 establishes the hard
direction of the dictionary for $G$. This is presumably what the original
draft \texttt{qwn\_atm\_1.txt} intended; making the induction explicit
(including the base case $r = 1$, the trivial subgroup cases handled by the
reductions, and a check that the smaller groups $G/H \cong \Z/p^{r-s}\Z$ and
$H \cong \Z/p^s\Z$ are both of order $< p^r$ when $0 < s < r$) is what is
required to make the present proof rigorous.
\end{remark}

\begin{remark}[Comparison with the direct Cartesian-inertia proof]
The proof of the same theorem in \texttt{claude\_atm\_1.tex} avoids
\Cref{prop:dictionary} entirely. It works directly with inertia groups in the
function-field tower $K_C \subseteq F \subseteq E$ inside its compositum with
$K_{C'}$, uses the Cartesian property of the master diagram to relate
$I_{v_E}(E/F) \cap N$ to the trivially-inertia subgroup
$I_{v_\star}(E\cdot K_{C'}/K_{C'})$ (which is trivial by (H2)), and then
deploys (H3) via the chain-of-subgroups property of cyclic $p$-groups to
conclude $I_{v_E}(E/F) = 1$.

That proof has three advantages over the present one:
\begin{itemize}
  \item It uses no local ramification theory: no Herbrand function, no upper
        numbering, no induction on $|G|$.
  \item It does not invoke the dictionary at all, so the circularity issue of
        \Cref{rem:circularity} does not arise.
  \item The role of (H3) is transparent --- it is the chain property of cyclic
        $p$-groups --- whereas the present proof loads (H3) into both
        \Cref{lem:herbrand-convexity} (via the chain of lower-numbering
        ramification groups) and the dictionary (via whatever proof of
        \Cref{prop:dictionary} one cites).
\end{itemize}
The present proof's only structural advantage is that
\Cref{lem:herbrand-convexity} is a clean, self-contained statement of local
ramification theory that may be of independent interest, and that the proof
factors cleanly through the (well-studied) local--global dictionary rather
than through an ad-hoc Galois argument inside a particular tower. For thesis
exposition the direct proof seems preferable; the Herbrand-route proof is
better seen as motivation for the lemma than as the primary proof.
\end{remark}

\printbibliography

\end{document}
